% !TeX root = ../main.tex
\section{Threat Model and Assumptions}
\label{sec:threat}

\textbf{Adversary.}
We consider a passive adversary (e.g., ISP or backbone observer) who records the
first 30 packets of each QUIC flow and extracts three metadata channels per packet:
normalized IPT, direction (client-to-server vs.\ server-to-client), and packet size.
The adversary does not decrypt payloads or QUIC headers protected by TLS~\cite{rfc9000}.
The adversary knows the feature extraction pipeline and has trained a supervised
multi-class classifier on historical labeled data from the same network environment.

\textbf{Defender.}
The defender operates an inline proxy that applies deterministic obfuscation to
outgoing (and optionally incoming) packets before they traverse the monitored link.
The defender cannot hide flow existence or volume at the coarsest level; the goal is
to reduce \emph{application-level} inferability.

\textbf{Out of scope.}
We do not model adaptive adversaries who retrain on obfuscated traces (see
Section~\ref{sec:discussion}), active attacks, or colluding endpoints. We evaluate a
single held-out test week; cross-network generalization is future work.

\textbf{Evaluation ethics.}
Experiments use published anonymized backbone traces under CESNET-QUIC22 terms of use;
no live user traffic is collected in this study.

%==============================================================================
